% Corrected French for printed p.126.
% Corrections: existe -> existent; diviseur -> diviseurs.
\SourcePage{126}

\LayerTarget{tohoku:section:1.3}
\noindent\textbf{1.3. Catégories additives.}
Une \textit{catégorie additive} est une catégorie $\mathbf{C}$ pour laquelle on s'est
donné, pour tout couple $(A,B)$ d'objets de $\mathbf{C}$, une loi de groupe abélien dans
$\mathrm{Hom}(A,B)$, de telle façon que la composition des morphismes soit une
opération bilinéaire.
\LayerTarget{tohoku:section:1.3:p1}
On suppose de plus que la somme et le produit de deux objets quelconques $A,B$
de $\mathbf{C}$ existent. Il suffit d'ailleurs de postuler l'existence de la somme
\textit{ou} du produit de $A$ et $B$, l'existence de l'autre s'en déduit alors
facilement et de plus $A+B$ est canoniquement isomorphe à $A\times B$.
(Supposant par exemple que $A\times B$ existe, on considère les morphismes
$A\to A\times B$ et $B\to A\times B$ dont les
composantes sont $(I_A,0)$ resp. $(0,I_B)$, et on vérifie qu'on obtient là une
représentation de $A\times B$ comme somme directe de $A$ et $B$). Enfin, on
postule l'existence d'un objet $A$ tel que $I_A=0$, on l'appelle un
objet nul ou \textit{zéro} de $\mathbf{C}$. Il revient au même de dire que
$\mathrm{Hom}(A,A)$ est réduit à zéro, ou encore que pour tout $B\in\mathbf{C}$,
$\mathrm{Hom}(A,B)$ (ou $\mathrm{Hom}(B,A)$) est réduit à zéro. Si $A$ et $A'$
sont deux objets nuls, alors il existe un isomorphisme unique de $A$ sur $A'$
(savoir l'unique élément $0$ de $\mathrm{Hom}(A,A')$ !), aussi on identifiera
tous les objets nuls de $\mathbf{C}$ à un seul, noté $0$ par abus d'écriture.

\LayerTarget{tohoku:section:1.3:p2}
La catégorie duale d'une catégorie additive est encore une catégorie additive.

\LayerTarget{tohoku:section:1.3:p3}
Soient $\mathbf{C}$ une catégorie additive, $u:A\to B$ un morphisme dans
$\mathbf{C}$. Pour que $u$ soit injectif (resp. surjectif) il faut et il suffit qu'il
n'existe pas de morphisme non nul, qui composé à droite (resp. à gauche) avec
$u$ donne zéro. On appelle \textit{noyau généralisé} de $u$ tout monomorphisme
$i:A'\to A$ tel que les morphismes $C\to A$ qui sont
diviseurs de zéro à droite de $u$ soient exactement ceux qui se factorisent en $C\to A'\mathrel{\mathop{\to}^{i}}A$; un tel monomorphisme est défini à équivalence près (cf.
N\textsuperscript{o}1), donc parmi les noyaux généralisés de $u$ (s'il en
existe) il y a exactement un qui soit un sous-truc de $A$ : on l'appelle
\textit{le noyau de $u$}, et on le note $\mathrm{Ker}\,u$. On définit
dualement le \textit{conoyau} de $u$ (qui est un truc quotient de $B$, s'il
existe), noté $\mathrm{Coker}\,u$. On appelle \textit{image} (resp.
\textit{coimage}) du morphisme $u$, le noyau de son conoyau (resp. le conoyau
de son noyau), s'il existe ; c'est donc un sous-truc de $B$ (resp. un truc
quotient de $A$)%
\LayerTarget{tohoku:section:1.3:footnote:1ter}\textsuperscript{1 ter)}
on les note $\mathrm{Im}\,u$ resp. $\mathrm{Coim}\,u$. Si $u$ admet une image
et une coimage, alors il existe un morphisme unique
$\bar u:\mathrm{Coim}\,u\to\mathrm{Im}\,u$ tel que $u$ soit
identique au composé $A\to\mathrm{Coim}\,u\to\mathrm{Im}\,u\to B$, où les morphismes extrêmes sont les morphismes canoniques.

\LayerTarget{tohoku:section:1.3:p4}
Un foncteur $F$ d'une catégorie additive $\mathbf{C}$ dans une autre $\mathbf{C}'$ est dit
\textit{foncteur additif} si pour deux morphismes
$u,v:A\to B$ dans $\mathbf{C}$, on a
$F(u+v)=F(u)+F(v)$. Définition analogue pour les multifoncteurs. Des foncteurs
composés de foncteurs additifs sont additifs. Si $F$ est un foncteur additif,
$F$ transforme une somme directe finie d'objets $A_i$ en la somme directe des
$F(A_i)$.

\SourceFootnote{1 ter}{\LayerTarget{tohoku:section:1.3:footnote:1ter:text}%
Une définition, plus naturelle à vrai dire, de
l'image de $u$, serait de prendre le plus petit sous-truc $B'$ de $B$ (s'il en
existe) tel que $u$ provienne d'un morphisme de $A$ dans $B'$. Cette définition
n'est équivalente à celle donnée dans le texte que dans le cas où $\mathbf{C}$ est une
catégorie \textit{abélienne} (cf. 1.4.)}

\TohokuContinuation{127}
